\pdfmapfile{+cm.map}
\pdfmapfile{+cmextra.map}
\documentclass[12pt,letterpaper]{article}
\usepackage[margin=0.65in,headheight=15pt,headsep=16pt,footskip=24pt]{geometry}
\usepackage{tikz}
\usepackage{fancyhdr}
\usetikzlibrary{arrows.meta}
\renewcommand{\familydefault}{\sfdefault}
\pagestyle{fancy}
\fancyhf{}
\fancyhead[L]{\fontsize{11}{13}\selectfont Week 77 / Persistent holes / Grades 4--5}
\fancyfoot[L]{\fontsize{9}{11}\selectfont Bellingham Math Circle / Week 77 / F77-45-v4}
\fancyfoot[R]{\fontsize{9}{11}\selectfont\thepage}
\renewcommand{\headrulewidth}{0pt}
\renewcommand{\footrulewidth}{0pt}
\hyphenpenalty=10000
\exhyphenpenalty=10000
\setlength{\parindent}{0pt}
\setlength{\parskip}{0pt}
\newcommand{\textat}[4]{\node[anchor=north west,inner sep=0pt,text width=#3cm,align=left,font=\fontsize{12}{16}\selectfont] at (#1,#2) {#4};}
\newcommand{\smallat}[4]{\node[anchor=north west,inner sep=0pt,text width=#3cm,align=left,font=\fontsize{11}{14}\selectfont] at (#1,#2) {#4};}
\newcommand{\lineat}[3]{\draw[black!28,line width=.4pt] (#1,#2)--(#3,#2);}
\newcommand{\token}[4]{\node[draw=black!60,rounded corners=1pt,minimum width=.83cm,minimum height=.5cm,inner sep=1pt,font=\fontsize{10}{12}\selectfont] (#4) at (#1,#2) {#3};}
\newcommand{\squareboard}[3]{%
 \begin{scope}[shift={(#1,#2)}]
 \coordinate (A) at (0,0); \coordinate (B) at (6.2,0); \coordinate (C) at (6.2,6.2); \coordinate (D) at (0,6.2);
 \draw[black!50,dash pattern=on 4pt off 4pt,line width=.9pt] (A)--(B)--(C)--(D)--cycle (A)--(C);
 \ifnum#3=1 \draw[line width=1.5pt] (A)--(B)--(C)--(D); \fi
 \foreach \v/\pos in {A/above left,B/above right,C/below right,D/below left}{\fill (\v) circle(2.2pt);\node[\pos,inner sep=4pt,font=\large] at (\v) {\v};}
 \end{scope}}
\newcommand{\pagebegin}{\begin{tikzpicture}[x=1cm,y=-1cm]\path[use as bounding box] (0,0) rectangle (18.2,22.15);}
\newcommand{\pageend}{\end{tikzpicture}}
% A problem file sets its prompt's height with \promptat{y} (centimetres from the top of the
% page's canvas); the prompt pattern in week.toml places the numbered prompt there.
\newcommand{\promptat}[1]{\def\prompty{#1}}\promptat{0}
\begin{document}
\pagebegin

\textat{0}{0}{18.2}{Keep track of a loop while you add edges and cover triangles with tiles. When do the filled tiles fill your saved loop?}
\smallat{0}{1.65}{18.2}{A loop follows edges back to its starting dot without repeating an edge. Save its edge list; keep that list unchanged as the board grows. Here the saved loop is the square's outside rim: AB, BC, CD, DA.}
\foreach \x/\stage/\caption in {0/0/{Build the rim},4.7/1/{Add the diagonal},9.4/2/{Fill one half},14.1/3/{Fill both halves}} {
 \begin{scope}[shift={(\x,4.25)}]
 \ifnum\stage>1 \fill[black!20] (0,0)--(3,0)--(3,3)--cycle;\fi
 \ifnum\stage>2 \fill[black!20] (0,0)--(3,3)--(0,3)--cycle;\fi
 \draw[line width=1.5pt] (0,0)--(3,0)--(3,3)--(0,3)--cycle;
 \ifnum\stage>0 \draw[line width=1pt] (0,0)--(3,3);\fi
 \foreach \a/\b/\v/\pos in {0/0/A/above left,3/0/B/above right,3/3/C/below right,0/3/D/below left}{\fill (\a,\b) circle(1.7pt);\node[\pos,inner sep=2pt,font=\small] at (\a,\b){\v};}
 \node[align=center,font=\fontsize{10}{12}\selectfont,text width=3.5cm] at (1.5,4.0){\caption};
 \end{scope}
}
\smallat{0}{9.0}{18.2}{One filled half does not fill the saved square rim. Both filled halves together do: their shared diagonal is inside, leaving the square rim as their outside boundary. We say the saved loop is \textbf{gone} when filled tiles have that loop as their combined boundary. The edges are still there!}
\smallat{0}{11.4}{18.2}{Use only printed edge places. Add pieces and keep them during a build. A partner checks that all three edges are in place before a triangle is filled. Dashed lines show edge places; solid lines are already present. AB and BA name the same edge.}
\smallat{0}{13.65}{18.2}{Here is a token test for a saved triangular loop. Copy its edge tokens, add the three edge tokens of the filled tile, then remove matching pairs. Nothing remains, so this loop is gone. Only test tokens are removed, never board edges or the saved list.}
\begin{scope}[shift={(.4,16.35)}]
 \fill[black!18] (0,1.4)--(1.8,1.4)--(.9,0)--cycle;
 \draw[line width=1pt] (0,1.4)--(1.8,1.4)--(.9,0)--cycle;
 \node[below left,inner sep=3pt] at (0,1.4) {X};\node[above,inner sep=3pt] at (.9,0) {Y};\node[below right,inner sep=3pt] at (1.8,1.4) {Z};
\end{scope}
\token{3.1}{16.9}{XY}{i1}\token{4.1}{16.9}{YZ}{i2}\token{5.1}{16.9}{XZ}{i3}
\smallat{2.7}{18.1}{3.4}{saved loop}
\draw[-{Stealth[length=2mm]},line width=.8pt] (5.8,17.35)--(7.2,17.35);
\token{8.1}{16.9}{XY}{a}\token{9.1}{16.9}{YZ}{b}\token{10.1}{16.9}{XZ}{c}
\token{8.1}{17.7}{XY}{d}\token{9.1}{17.7}{YZ}{e}\token{10.1}{17.7}{XZ}{f}
\foreach \n in {a,b,c,d,e,f}{\draw[line width=.8pt] (\n.south west)--(\n.north east);}
\smallat{7.4}{18.5}{4.5}{add filled XYZ;\newline cancel each pair}
\draw[-{Stealth[length=2mm]},line width=.8pt] (11,17.35)--(12.5,17.35);
\node[draw=black!40,minimum width=3.6cm,minimum height=1.1cm,font=\small] at (14.75,17.35){no tokens left};
\smallat{0}{20.3}{18.2}{For any saved loop, you may use any selection of the currently filled triangles in this test, each at most once. The loop is gone if some selection leaves no tokens; otherwise it survives.}

\pageend
\newpage
\pagebegin

\promptat{0}

\textat{0}{\prompty}{18.2}{\textbf{Problem 1:} Start with the dots only. Build a loop that can disappear after just one triangle is filled. Save its edge list. Start again and build a loop that survives either choice of a single filled triangle. Save that loop too. Test the two choices in separate builds.}

\squareboard{.65}{4.0}{0}
\smallat{8.2}{3.8}{9.7}{Loop that can disappear with one filling:}
\lineat{8.2}{5.1}{18.2}\lineat{8.2}{6.35}{18.2}
\smallat{8.2}{10.0}{9.7}{Loop that survives either single filling:}
\lineat{8.2}{11.3}{18.2}\lineat{8.2}{12.55}{18.2}
\lineat{0}{20.55}{18.2}\lineat{0}{21.95}{18.2}

\pageend
\newpage
\pagebegin

\promptat{0}

\textat{0}{\prompty}{18.2}{\textbf{Problem 2:} Save the first loop when it closes. Find every schedule in which that saved loop first disappears at stage 8. Explain why you have all of them.}

\squareboard{10.6}{3.4}{1}
\smallat{0.5}{3.1}{7.8}{Stage 0: AB, BC, CD}
\smallat{0.5}{4.45}{7.8}{Stage 2: one of DA and AC}
\smallat{0.5}{5.80}{7.8}{Stage 4: the other edge}
\smallat{0.5}{7.15}{7.8}{Stage 5: fill ABC or ACD}
\smallat{0.5}{8.50}{7.8}{Stage 8: fill the other triangle}
\lineat{0}{11.0}{18.2}\lineat{0}{12.5}{18.2}\lineat{0}{14.0}{18.2}\lineat{0}{15.5}{18.2}\lineat{0}{17.0}{18.2}\lineat{0}{18.5}{18.2}\lineat{0}{20.0}{18.2}\lineat{0}{21.5}{18.2}

\pageend
\newpage
\pagebegin

\promptat{0}

\textat{0}{\prompty}{18.2}{\textbf{Problem 3:} Start with the four solid edges. Add one dashed edge at stage 2 and the other at stage 4. Fill one triangle at stage 5 and the other at stage 8. Whichever dashed edge comes first, can the first loop be gone at stage 5? Answer the same question for the square in Problem 2. Save the first loop's edges each time, and record a schedule your partner can replay or explain why there is none.}

\begin{scope}[shift={(3.25,4.55)}]
 \coordinate (A) at (0,0);\coordinate (B) at (0,5.8);\coordinate (C) at (5.8,1.4);\coordinate (D) at (11.6,0);\coordinate (E) at (11.6,5.8);
 \draw[line width=1.5pt] (A)--(C)--(B) (D)--(C)--(E);
 \draw[black!50,dash pattern=on 4pt off 4pt,line width=.9pt] (A)--(B) (D)--(E);
 \foreach \v/\pos in {A/above left,B/below left,C/below,D/above right,E/below right}{\fill (\v) circle(2.2pt);\node[\pos,inner sep=4pt,font=\large] at (\v) {\v};}
\end{scope}
\smallat{0}{11.45}{8.7}{AB first:}
\lineat{0}{12.55}{8.7}\lineat{0}{13.85}{8.7}
\smallat{9.5}{11.45}{8.7}{DE first:}
\lineat{9.5}{12.55}{18.2}\lineat{9.5}{13.85}{18.2}
\smallat{0}{14.75}{8.7}{Square, AC first:}
\lineat{0}{15.85}{8.7}\lineat{0}{17.15}{8.7}
\smallat{9.5}{14.75}{8.7}{Square, DA first:}
\lineat{9.5}{15.85}{18.2}\lineat{9.5}{17.15}{18.2}
\textat{0}{18.0}{18.2}{Count each enclosed, unfilled region as one hole. Every build on these two boards has 0, 1, 2, 1, 0 holes after stages 0, 2, 4, 5, 8. Can those counts alone tell you when the first loop disappears?}
\lineat{0}{20.55}{18.2}\lineat{0}{21.85}{18.2}

\pageend
\newpage
\pagebegin

\promptat{0}

\textat{0}{\prompty}{18.2}{\textbf{Problem 4:} Build this schedule. Stage 4 is one step: place its edge, then its tile, and check only after both are in place. Does it ever have two holes at a finished stage? Move exactly one of those two additions to stage 3 or 5 so that it does. Find every legal change. Keep all other additions at their original stages.}

\squareboard{10.6}{5.3}{1}
\smallat{0.5}{5.15}{7.8}{Stage 0: AB, BC, CD}
\smallat{0.5}{6.55}{7.8}{Stage 2: DA}
\smallat{0.5}{7.95}{7.8}{Stage 4: AC and filled ABC}
\smallat{0.5}{9.35}{7.8}{Stage 8: filled ACD}
\lineat{0}{12.9}{18.2}\lineat{0}{14.2}{18.2}

\promptat{16.0}

\textat{0}{\prompty}{18.2}{\textbf{Problem 5:} Could a saved loop be gone at one stage and come back at a later stage, if no edge or filled triangle is removed? Explain why your answer works for every legal growing board.}

\lineat{0}{19.1}{18.2}\lineat{0}{20.45}{18.2}\lineat{0}{21.8}{18.2}

\pageend
\newpage
\pagebegin

\promptat{0}

\textat{0}{\prompty}{18.2}{\textbf{Problem 6:} Start with all eight edges and filled triangles ABO and BCO. Save two loops with different sets of edges. Both must survive now and first disappear at the last filling, whichever of the two unfilled triangles your partner fills first.}

\begin{scope}[shift={(9.5,4.15)}]
 \coordinate (A) at (0,0);\coordinate (B) at (8,0);\coordinate (C) at (8,8);\coordinate (D) at (0,8);\coordinate (O) at (4,4);
 \draw[line width=1.5pt] (A)--(B)--(C)--(D)--cycle (A)--(O)--(C) (B)--(O)--(D);
 \foreach \v/\pos in {A/above left,B/above right,C/below right,D/below left}{\fill (\v) circle(2.2pt);\node[\pos,inner sep=4pt,font=\large] at (\v) {\v};}
 \fill (O) circle(2.2pt);\node[above,fill=white,inner sep=2pt,font=\large] at (4,3.7) {O};
\end{scope}
\smallat{0}{4.15}{8.3}{First saved loop:}
\lineat{0}{5.35}{8.3}\lineat{0}{6.65}{8.3}
\smallat{0}{8.55}{8.3}{Second saved loop:}
\lineat{0}{9.75}{8.3}\lineat{0}{11.05}{8.3}
\textat{0}{13.65}{18.2}{Can you choose your pair so that either loop can disappear first on a fresh board with all eight edges and no filled triangles? Keep the two saved edge lists unchanged. Fill all four triangles, one at a time. Show a filling order for each outcome, or explain why it is impossible.}
\smallat{0}{17.4}{18.2}{First saved loop disappears first:}
\lineat{0}{18.65}{18.2}
\smallat{0}{20.0}{18.2}{Second saved loop disappears first:}
\lineat{0}{21.25}{18.2}

\pageend
\end{document}
